\section*{Chapitre \ref{ch:b3:gravity-geodesy} --- Pesanteur et forme de la Terre}

\begin{solution}{exo:b3:gravity-geodesy:1}
$b = a(1 - f) = 6378137\times0.99664719 = \qty{6356752.3}{m}$, et
$a - b = \qty{21384.7}{m}$, environ \qty{21.4}{km}. Clairaut~:
$f = (3J_2 + m)/2 = (0.00324789 + 0.00344979)/2 = \num{0.0033488}$, ou
$1/298.61$, contre $1/298.26$~: \qty{0.12}{\percent} de moins, soit la taille des
termes du second ordre négligés par la relation.
\end{solution}

\begin{solution}{exo:b3:gravity-geodesy:2}
$\sigma = \qty{6.0e13}{kg}/\qty{5.0e11}{m^2} = \qty{120}{kg.m^{-2}}$, donc
$h_w = 120/1000 = \qty{0.12}{m}$~: \qty{12}{cm} d'eau. Répartie sur
l'océan, $60/360 = \qty{0.17}{mm}$ du niveau de la mer.
\end{solution}

\begin{solution}{exo:b3:gravity-geodesy:3}
À l'équateur $\gamma = \qty{9.780327}{m.s^{-2}}$. À \ang{45},
$\sin^2\varphi = 0.5$ et $\sin^22\varphi = 1$~: $\gamma = 9.780327\times1.0026454
= \qty{9.806200}{m.s^{-2}}$. Aux pôles, $9.780327\times1.0053024 =
\qty{9.832186}{m.s^{-2}}$. La différence vaut \qty{0.051859}{m.s^{-2}}, soit
\qty{5186}{\milli\gal}.
\end{solution}

\begin{solution}{exo:b3:gravity-geodesy:4}
$D = ET_e^3/(12(1 - \nu^2)) = 10^{11}T_e^3/11.25$~: \qty{8.89e21}{N.m} pour
\qty{10}{km}, \qty{2.40e23}{N.m} pour \qty{30}{km}. Avec $(\rho_m - \rho_w)g =
\qty{22269}{N.m^{-3}}$, $\alpha = (4D/22269)^{1/4}$~: \qty{35.5}{km} et
\qty{81.0}{km}. La fosse flexurale s'achève à $3\pi\alpha/4$, \qty{84}{km} et \qty{191}{km}~;
le \omterm{def:b3:gravity-geodesy:alpha}{bombement périphérique} culmine à $\pi\alpha$, \qty{112}{km} et \qty{255}{km}.
\end{solution}

\begin{solution}{exo:b3:gravity-geodesy:5}
(a) $T = -\gamma N = -9.80\times50 = \qty{-490}{J.kg^{-1}}$. (b) Le géoïde s'incline
de $10/100\,000 = \num{1.0e-4}$ rad par rapport à l'ellipsoïde, et la verticale
(fil à plomb), normale au géoïde, s'incline du même angle par rapport à la normale de l'ellipsoïde~:
$\num{1.0e-4}\times206\,265 = 20.6''$.
\end{solution}

\begin{solution}{exo:b3:gravity-geodesy:6}
$H = J_2/(C/Ma^2) = \num{1.08263e-3}/0.3308 = \num{0.0032728}$, et
$1/H = 305.55$. $J_2$ est identique par hypothèse~; une Terre homogène aurait
$C/(Ma^2) = 0.4$ et donc $H = J_2/0.4 = \num{0.0027066}$~: la précession serait
plus lente. Mesurer à la fois $J_2$ et $H$ permet de mettre en évidence la
concentration de la masse vers le centre.
\end{solution}

\begin{solution}{exo:b3:gravity-geodesy:7}
$r = \qty{6878.137}{km}$~; $n = \sqrt{GM/r^3} = \qty{1.10678e-3}{rad.s^{-1}}$, une
période de $2\pi/n = \qty{94.6}{min}$. $\dot\Omega = -1.5\times\num{1.10678e-3}
\times\num{1.08263e-3}\times(6378.137/6878.137)^2\times\cos\ang{89} =
\qty{-2.70e-8}{rad.s^{-1}}$, soit \ang{-0.134} par jour, vers l'ouest. Une orbite presque polaire
survole toutes les latitudes, de sorte que la Terre entière est cartographiée~; elle précesse aussi
lentement, de sorte que les traces au sol se complètent régulièrement.
\end{solution}

\begin{solution}{exo:b3:gravity-geodesy:8}
$e^{-2\pi\times500/100} = e^{-31.4} = \num{2.3e-14}$~: \qty{2e-13}{\milli\gal},
rien. $e^{-2\pi\times500/1000} = e^{-3.14} = 0.0432$~:
\qty{0.43}{\milli\gal}.
\end{solution}

\begin{solution}{exo:b3:gravity-geodesy:9}
$3\pi\alpha/4 = \qty{120}{km}$ donne $\alpha = \qty{50.9}{km}$.
$D = (\rho_m - \rho_w)g\alpha^4/4 = 22269\times(50\,930)^4/4 =
\qty{3.75e22}{N.m}$, et $T_e = (11.25D/E)^{1/3} = \qty{16.2}{km}$. La plaque
franchit son niveau non fléchi à $\pi\alpha/2 = \qty{80}{km}$ de la fosse océanique.
\end{solution}

\begin{solution}{exo:b3:gravity-geodesy:10}
(a) Avec $u = x/\alpha$, $\mathrm d^4/\mathrm dx^4 = \alpha^{-4}\,\mathrm
d^4/\mathrm du^4$. En écrivant $\cos u + \sin u = \sqrt2\cos(u - \pi/4)$, la
fonction vaut $\sqrt2w_0\,\mathrm{Re}\,e^{(-1 + i)u - i\pi/4}$, et
$(-1 + i)^4 = -4$~: donc $w'''' = -4w/\alpha^4$. Puis $Dw'''' + (\rho_m -
\rho_i)gw = (-4D/\alpha^4 + (\rho_m - \rho_i)g)w = 0$, puisque $\alpha^4 =
4D/((\rho_m - \rho_i)g)$. (b) $w'' = -(2w_0/\alpha^2)e^{-u}(\cos u - \sin u)$, donc
$M = -Dw'' = (2Dw_0/\alpha^2)e^{-u}(\cos u - \sin u)$~: maximal en valeur absolue en
$u = 0$, nul là où $\tan u = 1$, soit $x = \pi\alpha/4$. (c) Sous la charge, la
plaque est concave vers le haut, sa surface supérieure en compression~; au-delà de $\pi\alpha/4$
elle est convexe vers le haut, sa surface supérieure étirée. Là, $M' \propto
-e^{-u}\cos u$ s'annule en $u = \pi/2$~: l'étirement est maximal en
$x = \pi\alpha/2$, sur le flanc externe de la fosse flexurale du côté du \omterm{def:b3:gravity-geodesy:alpha}{bombement périphérique}, où
$|M|$ vaut $e^{-\pi/2} = \qty{21}{\percent}$ de sa valeur sous la charge. C'est
là que des failles normales et des fissures s'ouvriraient au sommet de la plaque.
\end{solution}

\begin{solution}{exo:b3:gravity-geodesy:11}
$\sqrt{1 + \eta} = 2.5(1 - 1.5\times0.3308) = 1.2595$, donc $1 + \eta = 1.5863$
et $\tfrac52m/f = 2.5863$~: $f_h = 2.5\times0.00344979/2.5863 = \num{0.0033346}$,
$1/f_h = 299.9$. La différence $f - f_h = \num{1.82e-5}$~; multipliée par $a$, l'équateur
se situe à \qty{116}{m} au-dessus de la figure hydrostatique. La Terre n'est pas tout à fait
un fluide en équilibre~: son intérieur supporte des contraintes, dues à la résistance de
la lithosphère ou aux écoulements et aux contrastes de densité du manteau, qui maintiennent
cette forme un peu à l'écart de cet équilibre.
\end{solution}

\begin{solution}{exo:b3:gravity-geodesy:12}
Le nœud doit avancer de $2\pi/(365.242\times86\,400) =
\qty{1.991e-7}{rad.s^{-1}}$. Pour $r = \qty{7078.137}{km}$, $n =
\qty{1.0602e-3}{rad.s^{-1}}$ et $\tfrac32nJ_2(a/r)^2 = \num{1.398e-6}$, donc
$\cos i = -\num{1.991e-7}/\num{1.398e-6} = -0.1424$ et $i = \ang{98.2}$~: une
orbite légèrement rétrograde et héliosynchrone.
\end{solution}

\begin{solution}{pb:b3:gravity-geodesy:1}
\textbf{1.} $m = (\num{7.292115e-5})^2\times6378137^2\times6356752.3/
\num{3.986005e14} = \num{0.00344979}$.

\textbf{2.} $f = (3J_2 + m)/2 = \num{0.0033488}$, $1/f = 298.61$.

\textbf{3.} $\sqrt{1 + \eta} = 2.5(1 - 1.5\times0.3308) = 1.2595$,
$\tfrac52m/f_h = 2 + 0.5863$, donc $f_h = \num{0.0033346}$ et $1/f_h = 299.88$.

\textbf{4.} La valeur de Clairaut, $1/298.61$, se situe à moins de \qty{0.12}{\percent} de
la valeur observée $1/298.26$~; la valeur hydrostatique, $1/299.88$, est plus petite de
$\Delta f = \num{1.82e-5}$, et $a\,\Delta f = \qty{116}{m}$~: l'équateur se situe
\qty{116}{m} au-dessus de la figure hydrostatique.

\textbf{5.} $H = \num{1.08263e-3}/0.3308 = \num{0.0032728}$ ($1/305.6$).

\textbf{6.} $\tfrac12\times\qty{80}{km}\times\qty{6}{km} = \qty{240}{km^2} =
\qty{2.40e8}{m^2}$.

\textbf{7.} $V_0 = (2800 - 1030)\times9.81\times\num{2.40e8} =
\qty{4.17e12}{N.m^{-1}}$.

\textbf{8.} Airy~: $(\rho_m - \rho_v)r = (\rho_v - \rho_w)h$, donc $r =
6\times1770/500 = \qty{21.2}{km}$ sous la crête.

\textbf{9.} Le plancher océanique est fléchi bien au-delà du pied de la dorsale~: la
charge n'est pas compensée colonne par colonne mais par la flexion de la plaque sur une
région quatre fois plus large. La plaque possède une résistance~: compensation régionale.

\textbf{10.} $\alpha = 165/(3\pi/4) = \qty{70.0}{km}$.

\textbf{11.} $\alpha = 220/\pi = \qty{70.0}{km}$~: la même valeur, comme l'exige une
plaque élastique unique.

\textbf{12.} $D = (\rho_m - \rho_w)g\alpha^4/4 = 2270\times9.81\times(\num{7.00e4})^4/4
= \qty{1.34e23}{N.m}$.

\textbf{13.} $T_e = (12(1 - \nu^2)D/E)^{1/3} = (11.25\times\num{1.34e23}/
\num{1e11})^{1/3} = \qty{24.7}{km}$.

\textbf{14.} $w_0 = V_0\alpha^3/(8D) = \num{4.17e12}\times(\num{7.00e4})^3/
(8\times\num{1.34e23}) = \qty{1334}{m}$.

\textbf{15.} $w_0e^{-\pi} = 1334\times0.0432 = \qty{58}{m}$.

\textbf{16.} $2\times1334\times\num{7.00e4}\times2270\times9.81 =
\qty{4.16e12}{N.m^{-1}}$, égal à $V_0$ à l'arrondi de $D$ près.

\textbf{17.} \qty{24.7}{km} se situe dans la gamme de la dorsale hawaïenne. La
plaque s'enfonce de \qty{1.3}{km} sous l'axe au lieu de former une racine de
\qty{21}{km}~: la même poussée d'Archimède est répartie sur les \qty{330}{km} de la fosse.

\textbf{18.} $\lambda = 2\pi\times\qty{70.0}{km} = \qty{440}{km}$,
$k = 1/\alpha = \qty{1.43e-5}{m^{-1}}$.

\textbf{19.} $\sigma = (2900 - 1030)\times500 = \qty{9.35e5}{kg.m^{-2}}$, donc
$2\pi G\sigma = \qty{3.92e-4}{m.s^{-2}} = \qty{39.2}{\milli\gal}$ au fond de la mer~;
multiplié par $e^{-k\times\qty{5}{km}} = 0.931$, \qty{36.5}{\milli\gal} au niveau de la mer
(négatif là où la plaque est enfoncée).

\textbf{20.} $\sigma = (3300 - 2900)\times500 = \qty{2.0e5}{kg.m^{-2}}$~:
\qty{8.39}{\milli\gal}, multiplié par $e^{-k\times\qty{12}{km}} = 0.843$,
\qty{7.07}{\milli\gal}. Le total vaut \qty{43.6}{\milli\gal}.

\textbf{21.} $e^{-k\times\qty{400}{km}} = e^{-5.71} = 0.0033$~: \qty{0.14}{\milli\gal}.

\textbf{22.} $e^{-2\pi\times400/80} = e^{-31.4} = \num{2.3e-14}$~: le satellite
ne voit rien de la dorsale elle-même, seulement la flexion large, et encore
affaiblie trois cents fois~; l'anomalie propre de la dorsale doit être mesurée en mer.

\textbf{23.} $n \approx 2\pi\times6371/440 = 91$.

\textbf{24.} D'après les questions 10 à 13, l'\omterm{def:b3:gravity-geodesy:flexure}{épaisseur élastique effective} de la
plaque sous la dorsale d'Oru est
\textbf{\qty[text-series-to-math=true]{24.7}{km}}.
\end{solution}
